\chapter[\emj{🔤}~DP en Cadenas: LCS y Edit Distance]{\emj{🔤}~DP en Cadenas: LCS y Edit Distance}

\begin{dsaquees}

Comparar dos palabras o textos para ver qué tanto se parecen. Dos
preguntas clásicas:

\begin{itemize}
\item LCS (Longest Common Subsequence): ¿cuál es la secuencia de letras
      más larga que aparece en AMBAS, en el mismo orden (aunque no
      pegadas)?
\item Edit Distance: ¿cuántos cambios mínimos (insertar, borrar,
      sustituir una letra) necesito para convertir una palabra en la
      otra?
\end{itemize}

Ambas se resuelven con una tabla 2D donde cada celda compara un prefijo
de una cadena contra un prefijo de la otra.

\end{dsaquees}

\begin{dsanino}

Tienes dos collares de letras 📿 y quieres saber qué tanto se parecen.

Para el LCS: buscas las letras que puedes ir tachando en los dos collares
al mismo tiempo, en orden, formando la palabra escondida más larga que
comparten. Si las puntas coinciden, esa letra cuenta y avanzas en las
dos. Si no, pruebas quitar la punta de un collar o del otro y te quedas
con la mejor opción.

Para el Edit Distance: imagina que el corrector de tu celular 📱 quiere
cambiar una palabra por otra. Cuenta cuántos "arreglos" mínimos hace:
meter una letra, quitar una, o cambiar una por otra. Menos arreglos =
más parecidas. Escribes todo en una tabla para no repetir cuentas, y la
esquina te da la respuesta.

\end{dsanino}

\begin{dsaotraforma}

Ambos problemas usan una tabla \verb|dp[i][j]| que compara el prefijo de
largo $i$ de la cadena $A$ con el prefijo de largo $j$ de la cadena $B$.

\textbf{LCS} (longest common subsequence): si $A_i = B_j$ entonces
$dp[i][j] = dp[i-1][j-1] + 1$; si no, $dp[i][j] = \max(dp[i-1][j],\,
dp[i][j-1])$.

\textbf{Edit Distance} (distancia de Levenshtein): mínimo de inserciones,
borrados y sustituciones. Si las letras coinciden, no cuesta nada; si no,
$1 + \min$ de las tres operaciones.

Son la base de: diff de archivos, autocorrector, alineamiento de ADN y
detección de plagio.

\end{dsaotraforma}

\begin{dsadiagrama}
\begin{verbatim}
LCS de "ABCBDAB" y "BDCAB"   (respuesta: "BCAB", largo 4)

        ""  B  D  C  A  B
    ""   0  0  0  0  0  0
    A    0  0  0  0  1  1
    B    0  1  1  1  1  2
    C    0  1  1  2  2  2
    B    0  1  1  2  2  3
    D    0  1  2  2  2  3
    A    0  1  2  2  3  3
    B    0  1  2  2  3  4   <- LCS = 4

Diagonal +1 si las letras coinciden; si no, el máximo de izq/arriba.
\end{verbatim}
\end{dsadiagrama}

\begin{lstlisting}
CÓMO FUNCIONA
- Tabla (n+1) x (m+1); fila/columna 0 = cadena vacía.
- LCS: si A[i-1]==B[j-1] -> diagonal+1; si no -> max(arriba, izquierda).
- Edit Distance: si coinciden -> diagonal; si no -> 1 + min(insertar,
  borrar, sustituir) = 1 + min(dp[i][j-1], dp[i-1][j], dp[i-1][j-1]).

PSEUDOCÓDIGO (LCS)
──────────────────────────────────────────────────────────────

función lcs(A, B):
    dp = tabla (n+1) x (m+1) de ceros
    PARA i desde 1 hasta n:
        PARA j desde 1 hasta m:
            SI A[i-1] == B[j-1]:
                dp[i][j] = dp[i-1][j-1] + 1
            SINO:
                dp[i][j] = max(dp[i-1][j], dp[i][j-1])
    return dp[n][m]

CÓDIGO C++ (LCS)
──────────────────────────────────────────────────────────────

#include <string>
#include <vector>
#include <algorithm>
using namespace std;

int lcs(const string& A, const string& B) {
    int n = A.size(), m = B.size();
    vector<vector<int>> dp(n + 1, vector<int>(m + 1, 0));
    for (int i = 1; i <= n; i++)
        for (int j = 1; j <= m; j++)
            if (A[i-1] == B[j-1]) dp[i][j] = dp[i-1][j-1] + 1;
            else dp[i][j] = max(dp[i-1][j], dp[i][j-1]);
    return dp[n][m];
}

CÓDIGO C++ (Edit Distance / Levenshtein)
──────────────────────────────────────────────────────────────

int editDistance(const string& A, const string& B) {
    int n = A.size(), m = B.size();
    vector<vector<int>> dp(n + 1, vector<int>(m + 1));
    for (int i = 0; i <= n; i++) dp[i][0] = i;   // borrar todo A
    for (int j = 0; j <= m; j++) dp[0][j] = j;   // insertar todo B
    for (int i = 1; i <= n; i++)
        for (int j = 1; j <= m; j++)
            if (A[i-1] == B[j-1]) dp[i][j] = dp[i-1][j-1];  // sin costo
            else dp[i][j] = 1 + min({ dp[i][j-1],    // insertar
                                      dp[i-1][j],    // borrar
                                      dp[i-1][j-1] });// sustituir
    return dp[n][m];
}

COMPLEJIDAD (Big-O)
──────────────────────────────────────────────────────────────

  Problema        │ Tiempo     │ Espacio
  ────────────────┼────────────┼──────────────
  LCS             │ O(n·m)     │ O(n·m) o O(m) 1D
  Edit Distance   │ O(n·m)     │ O(n·m) o O(m) 1D

* Se puede reducir a O(m) espacio guardando solo la fila anterior.
\end{lstlisting}

\begin{dsaentrevistas}

✅ Subsequence vs Substring: LCS permite huecos (subsecuencia); "longest
   common substring" exige letras contiguas y usa otra recurrencia
   (reinicia a 0 cuando no coinciden). No los confundas.

💡 "Edit Distance" (LC 72) es de los DP 2D más preguntados. La clave son
   los tres vecinos: izquierda (insertar), arriba (borrar), diagonal
   (sustituir).

⚠️ Casos base primero: la fila 0 y la columna 0 no son ceros en Edit
   Distance: valen 0,1,2,3\ldots (convertir a/desde la cadena vacía).

🔑 Reconstruir la subsecuencia: recorre la tabla desde \verb|dp[n][m]|
   hacia atrás; cuando las letras coinciden, esa letra va en el LCS y te
   mueves en diagonal.

\end{dsaentrevistas}

\begin{dsaresumen}

• LCS: subsecuencia común más larga (permite huecos, respeta el orden).
• LCS: diagonal+1 si coinciden; si no, max(arriba, izquierda).
• Edit Distance: mínimos insert/borrar/sustituir para igualar cadenas.
• Edit Distance: 1 + min(izq, arriba, diagonal) si no coinciden.
• Ambos O(n·m); reducibles a O(m) espacio con una fila.
• Usos: diff, autocorrector, ADN, plagio.
════════════════════════════════════════════════════════════════

\end{dsaresumen}
